\section*{Capítulo \ref{ch:b3:electrode-kinetics} --- Cinética de electrodo y electroanálisis}

\begin{solution}{exo:b3:electrode-kinetics:1}
La pendiente catódica, $2.303RT/\alpha F = \qty{118}{mV}$, da $\alpha = 59.2/118 = 0.50$; $j_0 = \qty{2.0e-6}{A.cm^{-2}}$,
la ordenada en $\eta = 0$.
\end{solution}

\begin{solution}{exo:b3:electrode-kinetics:2}
$i_0 = \qty{2.0e-4}{A}$; $R_{\mathrm{ct}} = RT/(Fi_0) = 0.02569/\num{2.0e-4} = \qty{128}{\ohm}$.
\end{solution}

\begin{solution}{exo:b3:electrode-kinetics:3}
Para \ce{MgSO4}, $I = \frac12(4c + 4c) = 4c$. Con \qty{1.0}{mmol/L}, $I = \qty{4.0}{mol.m^{-3}}$ y $\kappa^{-1} =
\qty{0.96}{nm} \times \sqrt{100/4} = \qty{4.8}{nm}$; con \qty{0.10}{mol/L}, $I = \qty{400}{mol.m^{-3}}$, $\kappa^{-1} = \qty{0.48}{nm}$.
\end{solution}

\begin{solution}{exo:b3:electrode-kinetics:4}
$i_{\mathrm c} = C_{\mathrm{dl}}Av = \num{20e-6} \times 0.0707 \times 0.1 = \qty{0.14}{\micro A}$, menos del 1\,\% del pico. A
\qty{10}{V/s} vale \qty{14}{\micro A}, mientras que el pico solo crece hasta $19\sqrt{100} = \qty{190}{\micro A}$:
la corriente de carga crece como $v$, y la faradaica como $\sqrt v$.
\end{solution}

\begin{solution}{exo:b3:electrode-kinetics:5}
$D = \pi\bigl(i\sqrt t/nFAc^*\bigr)^2 = \pi(\num{12.2e-6}/(96\,485 \times 0.0707 \times \num{1.00e-6}))^2 =
\qty{1.0e-5}{cm^2.s^{-1}}$.
\end{solution}

\begin{solution}{exo:b3:electrode-kinetics:6}
$FAc^* = \qty{1.364e-2}{C.cm^{-1}}$, $F/RT = \qty{38.92}{V^{-1}}$. A \qty{50}{mV/s}: $\sqrt{38.92 \times 0.05 \times \num{7.0e-6}} =
\num{3.69e-3}$, $i_{\mathrm p} = 0.4463 \times \num{1.364e-2} \times \num{3.69e-3} = \qty{22.5}{\micro A}$; a \qty{500}{mV/s}, $\sqrt{10}$ veces
más, \qty{71}{\micro A}.
\end{solution}

\begin{solution}{exo:b3:electrode-kinetics:7}
(a) Reversible. (b) \omterm{def:b3:electrode-kinetics:reversibility}{Cuasirreversible}: la separación crece con la velocidad de barrido.
(c) Una etapa química sigue a la transferencia electrónica (EC): a velocidades bajas, el
producto se consume antes del barrido de vuelta; un barrido rápido se adelanta a la química.
\end{solution}

\begin{solution}{exo:b3:electrode-kinetics:8}
$K_{ij}a_j/a_i = \num{2e-4} \times 0.14/\num{4.0e-3} = 0.007$: el electrodo da un valor un 0.7\,\% demasiado alto.
\end{solution}

\begin{solution}{exo:b3:electrode-kinetics:9}
Mínimos cuadrados: $b = \qty{0.1093}{\micro A.L.\micro g^{-1}}$, $a = \qty{0.466}{\micro A}$, $s = \qty{0.011}{\micro A}$; $c_0 = a/b =
\qty{4.26}{\micro g.L^{-1}}$, con $u(c_0) = \frac sb\sqrt{\frac1n + \frac{\bar y^2}{b^2S_{xx}}} = \qty{0.12}{\micro g.L^{-1}}$
($n = 4$, $\bar y = 1.286$, $S_{xx} = 125$).
\end{solution}

\begin{solution}{exo:b3:electrode-kinetics:10}
$Q = 0.0500 \times 1800 = \qty{90.0}{C}$; $n(\ce{Cu}) = Q/2F = \qty{4.66e-4}{mol}$, \qty{29.6}{mg}. El
resultado solo depende de la carga, medida a partir de una corriente y un tiempo, y de la
constante de Faraday: no hace falta ningún patrón, siempre que la electrólisis sea
completa y su rendimiento faradaico sea del 100\,\%.
\end{solution}

\begin{solution}{exo:b3:electrode-kinetics:11}
$\sqrt{DFv/RT} = \sqrt{\num{1.0e-5} \times 38.92 \times 0.1} = \qty{6.24e-3}{cm.s^{-1}}$: $\Lambda = 1.6$, \omterm{def:b3:electrode-kinetics:reversibility}{cuasirreversible}.
A \qty{10}{V/s}, $\Lambda = 0.16$, todavía \omterm{def:b3:electrode-kinetics:reversibility}{cuasirreversible} pero más lejos del límite
reversible; la onda solo sería reversible por debajo de alrededor de \qty{1}{mV/s}.
\end{solution}

\begin{solution}{exo:b3:electrode-kinetics:12}
$s/b = \qty{0.0838}{mM}$ y $b^2S_{xx} = \qty{203.5}{\micro A^2}$. En el punto $\bar y$, la incertidumbre vale
$0.0838\sqrt{1 + 1/7} = \qty{0.090}{mM}$. Para una lectura de \qty{18.0}{\micro A}, vale
$0.0838\sqrt{1.143 + 83.0/203.5} = \qty{0.104}{mM}$. Con tres lecturas en $\bar y$, vale
$0.0838\sqrt{1/3 + 1/7} = \qty{0.058}{mM}$.
\end{solution}

\begin{solution}{pb:b3:electrode-kinetics:1}
\textbf{1.} $\ce{C6H12O6 -> C6H10O6 + 2 H+ + 2 e-}$ y $\ce{[Fe(CN)6]^3- + e- -> [Fe(CN)6]^4-}$;
reacción global:
\[
  \ce{C6H12O6 + 2 [Fe(CN)6]^3- -> C6H10O6 + 2 [Fe(CN)6]^4- + 2 H+} .
\]
En el electrodo, $\ce{[Fe(CN)6]^4- -> [Fe(CN)6]^3- + e-}$.
\textbf{2.} Dos.
\textbf{3.} El oxígeno disuelto en la sangre es escaso y variable, y su producto, el
peróxido de hidrógeno, solo se oxida a un potencial alto en el que reaccionan muchas otras
especies; el mediador está presente en un exceso conocido.
\textbf{4.} Para que su concentración en la superficie sea nula: la corriente queda entonces limitada
solo por la difusión e insensible a pequeños cambios de potencial.
\textbf{5.} La corriente disminuye como $t^{-1/2}$; las lecturas deben hacerse en el instante usado para el
calibrado.
\textbf{6.} Cottrell: $i\sqrt t = 42.0$, 41.4, 41.7, 41.6, 41.8: constante con una precisión del 1\,\%.
\textbf{7.} $D = \pi(\text{pendiente}/FAc)^2 = \pi(\num{41.8e-6}/(96\,485 \times 0.0300 \times \num{1.00e-5}))^2 = \qty{6.55e-6}{cm^2.s^{-1}}$.
\textbf{8.} $\sqrt{\pi \times \num{6.55e-6} \times 5} = \qty{0.010}{cm} = \qty{100}{\micro m}$: el empobrecimiento alcanza la parte superior de
la capa hacia los \qty{5}{s}; las lecturas posteriores quedarían por debajo de la ley de Cottrell.
\textbf{9.} La mitad: \qty{9.35}{\micro A}.
\textbf{10.} La corriente de fondo: oxidación de interferentes e impurezas,
carga de la doble capa, el hexacianoferrato(II) residual del propio mediador.
\textbf{11.} Sí: los residuos, como mucho de \qty{0.095}{\micro A}, se dispersan sin tendencia. Con mucha glucosa, el mediador, presente en
cantidad limitada, o la \omterm{def:b3:complex-kinetics:enzyme}{enzima} se saturan y la recta se curva.
\textbf{12.} $(7.10 - 0.356)/0.9189 = \qty{7.34}{mM}$.
\textbf{13.} $0.0838\sqrt{1 + 1/7 + (7.10 - 8.89)^2/203.5} = 0.0838 \times 1.076 = \qty{0.090}{mM}$.
\textbf{14.} El término en $1/m$ (una sola lectura); las lecturas repetidas, o varias tiras,
lo reducen.
\textbf{15.} $3 \times 0.077/0.919 = \qty{0.25}{mM}$.
\textbf{16.} $7.34 \times 180.16 = \qty{1322}{mg.L^{-1}} = \qty{132}{mg.dL^{-1}}$.
\textbf{17.} Añade su propia corriente de oxidación: el medidor da un valor demasiado alto.
\textbf{18.} Un \omterm{def:b3:electrode-kinetics:cv}{voltamperograma} de la tira sin glucosa muestra una onda anódica del
ascorbato allí donde se oxida el hexacianoferrato(II).
\textbf{19.} A un potencial más bajo se oxidan menos sustancias.
\textbf{20.} La corriente del blanco se resta de la lectura del electrodo con \omterm{def:b3:complex-kinetics:enzyme}{enzima}
antes de aplicar la recta de calibrado.
\textbf{21.} $D$ aumenta unos pocos por ciento por kelvin, y la corriente como $\sqrt D$: los medidores miden
la temperatura y la corrigen.
\textbf{22.} La viscosidad y la fracción de glóbulos rojos de la sangre cambian $D$ y la corriente;
unos patrones en una matriz similar anulan estos efectos.
\textbf{23.} Las variaciones de una tira a otra (área del electrodo, carga de \omterm{def:b3:complex-kinetics:enzyme}{enzima} y de mediador),
la temperatura y la composición de la sangre se suman a la incertidumbre propia del calibrado.
\textbf{24.} \textbf{$c(\text{glucosa}) = 7.34 \pm \qty{0.09}{mmol.L^{-1}}$} (incertidumbre típica), unos \qty{132}{mg.dL^{-1}}.
\end{solution}
